\section{اندازه‌گیری و آزمون نمونه}

\subsection{ضخامت، سفتی و جهت کشش}

ضخامت را در چند نقطۀ قابل‌برش بگیرید و کمینه و بیشینه را در برگۀ نمونه ثبت کنید. برای مقایسۀ دو پوست، از ابزار و فشار یکسان استفاده کنید. دهانۀ کولیس را به‌آرامی ببندید تا چرم نرم فشرده نشود. برای گزارش قابل‌مقایسۀ آزمایشگاهی، روش تعیین ضخامت \lr{ISO 2589:2016} را مبنا قرار دهید \cite{iso2589}.

در فروش چرم ممکن است ضخامت را با \lr{oz} بیان کنند. در این قرارداد هر \lr{1 oz} تقریباً معادل \lr{0.4}~میلی‌متر است و بازۀ \lr{3--4 oz} حدود \lr{1.2--1.6}~میلی‌متر می‌شود. ضخامت واقعی و دامنۀ تغییر آن را برای همان سری کالا بپرسید \cite{tandy}.

سفتی در خمش، رفتار چرم در برابر تا کردن است. دو تکه با ضخامت برابر می‌توانند سفتی متفاوتی داشته باشند. تکه‌هایی هم‌اندازه را در جهت برش آینده آماده کنید، روی یک استوانه با شعاع یکسان خم کنید و نیروی تقریبی لازم، برگشت شکل و آثار خط تا را مقایسه کنید. برای کشش نیز دو تکه در جهت‌های عمود بر هم بگیرید و با نیروی دستی مشابه مقایسه کنید. این دو آزمون برای رتبه‌بندی نمونه‌های همان کارگاه‌اند. برای مقایسۀ کمّی یا قراردادی، روش آزمایش و شرایط نمونه باید مشخص باشند.

\subsection{پرداخت، انتقال رنگ و لبه}

\begin{enumerate}
  \item \textbf{تا.} رویه را در سمت بیرونیِ خم واقعی قرار دهید. اثر سفیدشدن موقت را از ترک پایدار، پوسته‌شدن پوشش و چروک عمیق جدا کنید. برای نمونه‌های پرزدار، جهت پرز را قبل و بعد از خم مقایسه کنید.
  \item \textbf{انتقال رنگ.} پارچۀ سفید را روی تکۀ دورریز بکشید. برای غربال اولیۀ دو نمونه، تعداد حرکت و فشار را یکسان نگه دارید. در صورت نیاز، آزمون نم‌دار را روی تکۀ دیگری انجام دهید. این مقایسه جای آزمون استاندارد ثبات رنگ در برابر مالش رفت‌وبرگشتی را نمی‌گیرد \cite{iso11640}.
  \item \textbf{لبه.} دو تکه را با ترتیب لایه‌ها، چسب و نخ واقعی پروژه آماده کنید. یک طرف را صیقل دهید و طرف دیگر را با پرداخت لبۀ مدنظر آزمایش کنید. پس از خشک‌شدن، یکنواختی، بلندشدن الیاف و تغییرِ پس از چند بار خم را بررسی کنید. چرم پرزدار یا نرم ممکن است با لبۀ تاخورده یا پرداخت پوششی بهتر با طرح سازگار شود.
  \item \textbf{دوخت.} نزدیک‌ترین فاصلۀ درز تا لبه را روی تکۀ آزمایشی اجرا کنید. اگر سوراخ‌ها به هم نزدیک‌اند و پل میان آن‌ها تغییر شکل می‌دهد، فاصلۀ سوراخ‌ها، فاصلۀ درز از لبه، قطر سوراخ و ضخامت لایه‌ها را با هم بازبینی کنید. استحکام اتصال یراق را با درز مستقیم یکی فرض نکنید.
\end{enumerate}

\section{انتخاب برای جاکارتی، کیف و بند}

از قطورترین نقطۀ مونتاژ شروع کنید. ضخامت لایه‌ها را در همان مقطع جمع بزنید، سپس مجموع واقعی را پس از چسب و دوخت اندازه بگیرید. اگر مجموع در خط تا یا زیر یراق زیاد شد، تغییر را در قطعات هم‌پوشان بررسی کنید. نازک‌کردن موضعی باید روی نمونۀ دوخته‌شده آزموده شود تا از خط تا یا محل اتصال ضعیف عبور نکند.

\begin{table}[htbp]
  \centering
  \caption{بازه‌های پیشنهادی آغاز آزمون برای سه طرح، به میلی‌متر}
  \label{tab:projects}
  \begin{tabular}{@{}p{23mm}p{60mm}p{63mm}@{}}
    \toprule
    محصول & ضخامت پیشنهادی هر قطعه & ناحیۀ تعیین‌کننده \\
    \midrule
    جاکارتی & بدنه \lr{1.0--1.4} و جیب‌ها \lr{0.6--0.9} & گوشۀ مشترک جیب‌ها و دهانۀ جای کارت. ورود و بیرون‌آمدن کارت را با ظرفیت واقعی امتحان کنید. \\
    کیف کوچک & رویه \lr{1.2--1.8} و آستر چرمی \lr{0.6--0.9} & درز گوشه، محل برگشت لبه و اتصال بند. سفتی رویه را پس از افزودن آستر بسنجید. \\
    بند تک‌لایه & \lr{3.0--4.0} برای شروع بررسی بند کمربند & بخش برگشتی در سگک و فاصلۀ سوراخ تا لبه. طول مفید، کشش در راستای بند و جاگیری در یراق را بررسی کنید. \\
    \bottomrule
  \end{tabular}
\end{table}

بازه‌های جدول~\ref{tab:projects} پیشنهادهای آزمون این جزوه برای ابعاد و ساخت‌های متعارف‌اند. آن‌ها برآورد طراحی هستند و با الگو، سفتی چرم، شکل اتصال و بار کاری تغییر می‌کنند. جدول ضخامت فروشنده نیز فقط نقطۀ شروع انتخاب است \cite{tandy}.

\subsection{جاکارتی با دو جیب روی هم}

فرض کنید بدنۀ جاکارتی \lr{1.2}~میلی‌متر و هر جیب \lr{0.8}~میلی‌متر ضخامت دارد. در گوشۀ مشترک سه قطعه، جمع اسمی لایه‌ها \lr{2.8}~میلی‌متر است. در نقطه‌ای که فقط بدنه و یک جیب روی هم می‌افتند، جمع به \lr{2.0}~میلی‌متر می‌رسد. به همین دلیل ممکن است یک شیوۀ سوراخ‌کاری و پرداخت، در امتداد لبه نتیجۀ یکنواختی ندهد.

ابتدا گوشۀ مشترک را نمونه‌سازی کنید. اگر لبه بیش از حد ضخیم شد، ترتیب هم‌پوشانی را اصلاح کنید یا پشت لایۀ مناسب را در لبۀ اتصال نازک کنید. پس از اجرا، دهانۀ جیب را با کارت پر کنید و جا افتادن کارت، شکل گوشه و دوام سوراخ‌های درز را دوباره ببینید.

\subsection{کیف نرم با آستر}

در کیف نرم، رویه و آستر را به‌صورت جداگانه و سپس به‌شکل مونتاژشده خم کنید. آستر ممکن است افت آزاد بدنۀ کیف را کم کند یا در گوشه کمی جمع شود. در محل اتصال بند، نمونه‌ای با تعداد واقعی لایه‌ها بسازید و تغییر شکل دور سوراخ یا یراق را ببینید. در صورت نیاز، تقویت را در ناحیۀ اتصال متمرکز کنید تا بقیۀ بدنه رفتار موردنظر را حفظ کند.

\subsection{بند کمربند با سگک}

برای بند، یکنواختی ضخامت و طول قابل‌برش را در کل نوار ارزیابی کنید. اگر بند \lr{3.5}~میلی‌متر باشد، دو لایه روی هم در بخش برگشتی سگک پیش از نازک‌کاری و مونتاژ حدود \lr{7.0}~میلی‌متر می‌شوند. ابتدا نتیجۀ تا را با خودِ سگک و یراقِ انتخابی بررسی کنید. برای بندی که بار مشخصی را تحمل می‌کند، ابعاد کمربند راهنمای تعیین ظرفیت بار نیست. برای چنین طرحی، داده یا آزمون متناسب با بار، اتصال و شیوۀ کاربرد لازم است.

\section{چک‌لیست خرید}

\begin{enumerate}
  \item نام و کد محصول، گونه، لایه، نوع دباغی و پرداخت را از فروشنده بخواهید. اصطلاح‌های تجاری مبهم را در مشخصات دقیق کنید.
  \item بازۀ ضخامت را به میلی‌متر بپرسید و سفتی موردنظر را با کاربرد شرح دهید. مثلاً رویۀ کیف می‌تواند نرم باشد و قطعۀ اتصال بند رفتار سفت‌تری بخواهد.
  \item رویه، پشت و مقطع را در نور ثابت بررسی کنید. محل عیب‌ها و سطح یا طول مفید برای برش را با بزرگ‌ترین قطعۀ الگو مقایسه کنید.
  \item برای سفارش مجدد، نمونۀ مرجع و کد سری را نگه دارید. امکان اختلاف رنگ یا پرداخت میان سری‌ها را در قطعات هم‌نما در نظر بگیرید.
  \item آزمون خم، مالش، لبه و دوخت را متناسب با طرح انجام دهید. نتیجه را کنار مشخصات، تاریخ و کاربرد نمونه یادداشت کنید.
  \item پیش از برش اصلی، دست‌کم یک مقطع چندلایه و یک گوشۀ واقعی را با الگو، نخ، چسب و یراق نهایی بسازید.
\end{enumerate}

در گزارش آزمون، نام نمونه و کد سری، محل اندازه‌گیری، بازۀ ضخامت، جهت برش، رفتار خم، وضعیت لبه پس از پرداخت و تصمیم برای هر قطعۀ الگو را کنار هم ثبت کنید. اگر چند نمونه به آزمون راه یافتند، معیار قبولی هر قطعه را پیش از مقایسه تعیین کنید.
